\chapter{Related Works}

This chapter reviews previous academic research and open-source systems that address credit card fraud detection. It analyzes traditional models, gradient boosting frameworks, recurrent networks, class imbalance techniques, and model explainability, establishing how FraudNet addresses current limitations.

\section{Review of Existing Works}

Over the past two decades, credit card fraud detection has shifted from manual verification and rule-based systems to automated statistical modeling and machine learning. Traditional baseline systems, such as those analyzed by Dal Pozzolo et al. \cite{pozzolo2015calibrating}, relied heavily on linear classifiers like Logistic Regression and basic Decision Trees. While computationally efficient, these models struggle to capture non-linear feature interactions and are highly sensitive to outliers, limiting their predictive accuracy in high-dimensional tabular datasets.

To improve predictive performance, tree-based ensemble methods were introduced. Breiman \cite{breiman2001random} developed the Random Forest (\gls{rf}) algorithm, which reduces variance by averaging predictions across multiple randomized decision trees. Ensembles like \gls{rf} are robust to overfitting and handle non-linear feature splits well. However, they lack direct sequential memory, processing each credit card transaction in isolation. Consequently, they fail to detect patterns of behavioral changes over time.

More recently, gradient boosting frameworks have dominated tabular classification challenges. Chen and Guestrin \cite{chen2016xgboost} introduced XGBoost, an optimized gradient boosting system that uses parallel tree boosting and regularization to achieve state-of-the-art accuracy. Similarly, Ke et al. \cite{ke2017lightgbm} developed LightGBM, which uses gradient-based one-side sampling and exclusive feature bundling to speed up training on large-scale datasets. While both models perform exceptionally well on tabular datasets, they still analyze transactions as independent events.

To capture temporal patterns, researchers began applying sequence-based models to fraud detection. Hochreiter and Schmidhuber \cite{hochreiter1997lstm} introduced Long Short-Term Memory (\gls{lstm}) networks, a class of recurrent networks designed to learn long-term dependencies. In fraud detection pipelines, \gls{lstm} and \gls{rnn} architectures process sequences of historical transactions to learn customer-specific spending behaviors. However, sequence models are computationally expensive, require substantial history for training, and do not provide clear explanations for their predictions.


\section{Addressing Class Imbalance}

One of the most persistent issues in credit card fraud detection is extreme class imbalance. Chawla et al. \cite{chawla2002smote} developed the Synthetic Minority Over-sampling Technique (\gls{smote}), which generates synthetic examples along the line segments joining k-nearest neighbors in the minority class. While widely used, applying \gls{smote} globally across a dataset before partitioning it into train and test sets leads to information leakage, artificially inflating evaluation metrics. FraudNet addresses this by applying \gls{smote} strictly within stratified, group-aware training folds.

\section{Explainable AI and Interpretability}

As fraud detection models grow more complex, interpretability becomes essential for validation and regulatory compliance. Lundberg and Lee \cite{lundberg2017unified} proposed SHAP (Shapley Additive exPlanations), a game-theoretic approach that explains individual predictions by calculating the contribution of each feature to the final score. Integrating \gls{shap} into real-time scoring systems allows risk analysts to verify why a transaction was flagged as fraudulent, resolving the "black box" limitation of deep learning and ensemble models.

